% ตัวอย่างเฉลยพาร์ทคอมพิวเตอร์ที่ถูกต้อง 1 ข้อ — แม่แบบของ "ครบ 6 ส่วน"
% ห้าม \input ไฟล์นี้เข้าไปในเฉลยจริง เปิดอ่านอย่างเดียว
%
% รูปแบบถอดมาจาก ~/Documents/POSN.Computer/ข้อสอบเทียม/ตัวอย่างไฟล์เฉลย คอม.pdf
% (30 ข้อ ทุกข้อมีครบ 6 ส่วน และมีอย่างน้อย 3 ขั้น)
%
% ไฟล์เฉลยพาร์ทคอมใส่ \renewcommand{\anskeyword}{เฉลย} ไว้หลัง \begin{document}
% หัวข้อจึงขึ้นว่า "N. เฉลย ค. 21" ตามไฟล์อ้างอิง
%
% โจทย์: โปรแกรมข้างล่างพิมพ์ค่าอะไร
%     n = 6
%     c = 0
%     i = 1
%     while i <= n:
%         j = i
%         while j <= n:
%             c = c + 1
%             j = j + 1
%         i = i + 1
%     print(c)
%
% ตัวเลือกทั้งสี่มาจาก verify.py:
%   R[7] = 21
%   D[7] = {'ลืมนับแนวทแยง (j เริ่มที่ i+1)': 15,
%           'คิดว่าเป็นครึ่งหนึ่งของตารางเต็มพอดี': 18,
%           'คิดว่า j เริ่มที่ 1 ทุกรอบ (ตารางเต็ม)': 36}

\Needspace{4\baselineskip}
\item
\ansline{ค.\ $21$}{}

\idea{ลูปชั้นในเริ่มที่ \code{j = i} ไม่ใช่ \code{j = 1} จำนวนรอบชั้นในจึง
\textbf{ลดลงทีละหนึ่ง} ทุกครั้งที่ \code{i} เพิ่ม ข้อนี้ $n=6$ ซึ่งเล็กพอ
วิธีที่เร็วและตรวจง่ายที่สุดคือ \textbf{ไล่ตาราง} ไม่ต้องนึกสูตรออกก็ทำได้}

\howto
\stepa{1}{ดูว่ารอบชั้นในทำกี่ครั้งต่อ \code{i} หนึ่งค่า \code{j} เริ่มที่ \code{i}
          และวิ่งจนถึง $n=6$ จึงทำได้ $6-i+1$ ครั้ง}
\stepa{2}{ไล่ตารางให้ครบทุกค่าของ \code{i} (มีแค่ 6 แถว เขียนออกมาดูเลย)
  \par\vspace{3pt}\centering
  \begin{tabular}{|c|c|c|}\hline
    \code{i} & \code{j} วิ่งจาก--ถึง & จำนวนรอบชั้นใน \\\hline
    1 & 1--6 & 6 \\
    2 & 2--6 & 5 \\
    3 & 3--6 & 4 \\
    4 & 4--6 & 3 \\
    5 & 5--6 & 2 \\
    6 & 6--6 & 1 \\\hline
  \end{tabular}\par\vspace{3pt}}
\stepa{3}{\code{c} เพิ่มหนึ่งทุกรอบชั้นใน ค่าสุดท้ายจึงเป็นผลบวกของคอลัมน์ขวา
          \[ c = 6+5+4+3+2+1 = 21 \]}
\stepa{4}{ตรวจซ้ำด้วยสูตรผลบวก $1$ ถึง $n$ ซึ่งตรงกับรูปสามเหลี่ยมที่ไล่ได้
          \[ \frac{n(n+1)}{2}=\frac{6\times7}{2}=21 \quad\checkmark \]}
\ansfinal{\code{21}}

\whywrong{\quickck{ลูปซ้อนที่ชั้นในเริ่มที่ \code{i} นับได้ \textbf{ครึ่งบนของตาราง
รวมแนวทแยง} จึงต้องมากกว่าครึ่งหนึ่งของ $6^2=36$ คือมากกว่า $18$ แต่ยังน้อยกว่า $36$
เหลือตัวเลือกเดียวทันที}
ก. $15$ คือลืมนับแนวทแยง (คิดว่า \code{j} เริ่มที่ \code{i+1})
ข. $18$ คือคิดว่าเป็นครึ่งหนึ่งของตารางเต็มพอดี
ง. $36$ คือคิดว่า \code{j} เริ่มที่ \code{1} ทุกรอบ จึงได้ตารางเต็ม $6\times6$}

\pitfall{อ่านลูปซ้อนแล้วคูณจำนวนรอบสองชั้นเข้าด้วยกันทันที ($6\times6$) ทั้งที่ขอบเขต
ของชั้นในขึ้นกับตัวนับชั้นนอก \textbf{ถ้าขอบเขตชั้นในมีตัวแปรของชั้นนอกอยู่ ห้ามคูณ
ให้ไล่ตารางแทน}}

% สิ่งที่ตัวอย่างนี้ทำถูกและต้องทำตามทุกข้อ
%
% 1. ครบ 6 ส่วนตามลำดับ: \ansline \idea \howto \stepa \ansfinal \whywrong \pitfall
% 2. 4 ขั้น (ขั้นต่ำคือ 3) โดยขั้นที่ 4 เป็นการตรวจซ้ำอีกทาง ไม่ใช่การยืดให้ครบโควตา
% 3. เลือก "ไล่ตาราง" เพราะ n=6 เล็กพอ — เห็นภาพและตรวจง่ายกว่าการนึกสูตรออกก่อน
%    ส่วนสูตรใส่ไว้เป็นตัวตรวจซ้ำท้ายข้อ ตามลำดับความสำคัญใน docs/ANSWER-KEY.md
% 4. โค้ดทุกชิ้นอยู่ใน \code{} — ต้องใช้ฟอนต์ monospace ไม่ใช่ตัวหนังสือไทยธรรมดา
% 5. \quickck ตัดตัวเลือกด้วยการประมาณขอบเขต ไม่ใช่การไล่ตารางซ้ำรอบสอง
% 6. ทุกตัวลวงมีที่มาจาก D[] และคำอธิบายคือชื่อความผิดที่รันเป็นโค้ดแล้วได้ตัวเลขนั้นจริง
